
\section{Anticipated questions}\label{app:faq}
\paragraph{Why is there no wet-lab validation?}
Scope and access: this is a computational project, and every candidate is
therefore stated as a falsifiable prediction with a named assay
(Appendix~\ref{app:cand-full}) rather than a discovery claim. The MIC
protocol is specified so an outside lab could run it directly.

\paragraph{Why not retrain on the reviewed-only subgroup?}
The reviewed-positive subgroup has 239 positives against 2{,}925 reviewed
negatives (Appendix~\ref{app:strata}). A model retrained on it would be
twelve times outnumbered at the class level and tiny in absolute terms;
its intervals would swallow any effect. The correct fix is a fresh,
length-matched, reviewed-only corpus at the original scale, which is a
data effort, not a re-run, and is recorded as the top future-work item.

\paragraph{Why is there no significance test against the published
Feb2020 row?}
A paired test needs either the published fold assignments or the
published per-sequence outputs; neither is available
(Appendix~\ref{app:sources}). With only the published means, any
``significance'' would be computed against an assumed variance and would
manufacture confidence we do not have. We report the descriptive gap and
its direction instead, and Table~\ref{tab:feb2020} shows fold-to-fold
variation larger than the gap.

\paragraph{Why is the frozen ensemble used for the screen after the
training-set confound was found?}
The confound was discovered after the screen ran, and retraining a clean
ensemble is the future-work item above. The screen results are retained
because the ranking signal (Section~\ref{sec:labelconfound}) plausibly
survives the confound while the calibration does not, and because
withdrawing them would hide exactly the audit trail this report exists to
document. Every consumer-facing number carries the caveat.

\paragraph{Why do the GNN results stay in the paper if they are the
weakest models?}
Because they are measured, honest, and informative: the capacity and
pooling analysis (Appendix~\ref{app:arch}) explains the gap, the GNN adds
ensemble diversity on the Veltri stack, and deleting weak results is how
benchmark literature becomes unreliable.
